\documentclass[11pt,a4paper]{article}
\usepackage[margin=25mm]{geometry}
\usepackage{fontspec}
\usepackage{xeCJK}
\setmainfont{Times New Roman}
\setCJKmainfont{Microsoft JhengHei}
\usepackage{amsmath,amssymb,amsthm,mathtools}
\usepackage{booktabs,longtable,array}
\usepackage{enumitem}
\usepackage{microtype}
\usepackage[numbers,sort&compress]{natbib}
\usepackage[hidelinks]{hyperref}
\usepackage{xurl}

\newtheorem{definition}{Definition}[section]
\newtheorem{proposition}[definition]{Proposition}
\newtheorem{conjecture}[definition]{Conjecture}
\newcommand{\formalized}{\textsc{Formalized}}
\newcommand{\derived}{\textsc{Derived}}
\newcommand{\conjectured}{\textsc{Conjecture}}
\newcommand{\powerset}{\mathcal P}
\newcommand{\true}{\mathsf{true}}
\newcommand{\false}{\mathsf{false}}

\title{Compositional Assurance for Legal Reasoning:\\
Request Identity, Horn Closure, Argumentation, Temporal Scope, and Human Authority}
\author{Laupinco}
\date{September 2026}

\begin{document}
\maketitle

\begin{abstract}
This paper develops a formula-centered, machine-checked architecture for legal reasoning. Sixteen Unified Legal Model modules formalize request identity, fail-closed outcomes, typed transitions, proof obligations, finite execution, Horn least closure, structured arguments, attacks and Dung semantics, branch-sensitive queries, procedure-aware adjudication, exact dimension-indexed calculation, conservative trust composition, add-only incrementality, conditional fixed-point convergence, temporal applicability, authority receipts, and taint non-interference. Each claim is marked \formalized{}, \derived{}, or \conjectured{}. The proved contribution is a family of local composition constraints: under their stated premises, transformations preserve request identity, branch-sensitive results do not compose across incompatible branches, obligation and trust joins do not hide weaker inputs, temporal guards exclude future evidence, and repeated machine agreement does not create higher authority. A separate release policy requires evidence that these constraints and their executable gates hold for one named source subject. The paper does not claim one global Lean theorem for release, substantive legal correctness, complete runtime refinement, or universal jurisdictional coverage.
\end{abstract}

\begin{center}\textbf{中文摘要}\end{center}

本文提出一套以公式为中心、由机器核验的法律推理架构。十六个统一法律模型模块形式化请求身份、失败闭合结果、类型化转移、证明义务、有限执行、Horn 最小闭包、结构化论证、攻击与 Dung 语义、分支查询、程序性裁判、维度索引精确计算、保守信任组合、只增量更新、条件式不动点收敛、时间适用性、权限收据与污点不干扰。所有结论分别标记为 \formalized{}、\derived{} 或 \conjectured{}。已证明的贡献是一组局部组合约束：在各自明示前提下，转换保持请求身份，分支结果不能跨不相容分支组合，义务与信任聚合不能隐藏较弱输入，时间门禁排除未来证据，重复机器意见不能产生更高权限。另由独立发布策略要求这些约束及其可执行门禁都绑定同一源码对象。本文不主张存在统一的 Lean 发布定理，也不主张实体法律正确性、完整运行时精化或跨法域普遍覆盖。

\paragraph{Keywords.} computational law; formal methods; Lean 4; argumentation; non-monotonic reasoning; assurance cases.

\section{Introduction}

This section has been rewritten as a release-bounded statement for the public repository. The section explains the specification boundary only. It does not assert full runtime correctness or private workflow validity.

The controlling sources are the Lean files in \texttt{proofs/lean/juris\_lean/JurisLean/}, the canonical specification files in \texttt{docs/spec/}, and the machine-readable manifests in \texttt{docs/formal-release/}. Generated reports and narrative papers are explanatory material.

For the current public boundary, the relevant slices are contract breach, license, permission, and priority. Each slice must stay within the canonical legal types, the minimal DDL core, the Horn-to-AAF contract, and the certificate/checker boundary. Missing or stale evidence is fail-closed.

The 2026-07-03 juris-calculus absorption of transferable LSC ideas is a runtime governance update. Fact trust envelopes, downgrade protocols, provenance fields, taint markers, output firewalls, cross-module IO declarations, conflict certificates, and review packets remain engineering metadata unless a separate source-level change modifies the formal specification.

\section{Preliminaries}
\label{sec:preliminaries}

\subsection{Requests, observations, and outcomes}

A request key (k) contains a declared case scope and a run scope. Its minimal well-formedness condition is
\begin{equation}
\operatorname{WF}(k)\iff
k.\operatorname{runScope.caseScope}=k.\operatorname{caseScope}.
\label{eq:request-wf}
\end{equation}
This \formalized{} equality prevents accidental cross-case substitution after normalization. It does not authenticate a real-world case identifier. For representation change $f:A\to B$, observations $o_A:A\to O$ and $o_B:B\to O$ define
\begin{equation}
\operatorname{Preserves}(o_A,o_B,f)
\iff\forall x,\ o_B(f(x))=o_A(x).
\label{eq:observation-preservation}
\end{equation}
Preservation composes: if $f$ preserves from $A$ to $B$ and $g$ preserves from $B$ to $C$, then $g\circ f$ preserves from $A$ to $C$. This is \formalized{}. The observation function remains a specification choice; omitting a legally material coordinate from $o$ cannot be repaired by the theorem.

Partial information is represented by a three-way sum:
\begin{equation}
\operatorname{Outcome}(\alpha)=
\operatorname{Complete}(\alpha)+
\operatorname{Partial}(\alpha,O\neq\varnothing)+
\operatorname{Failure}(e).
\label{eq:outcome}
\end{equation}
The witness $O\neq\varnothing$ ensures that partiality records at least one unresolved obligation. A generic map cannot promote failure:
\begin{equation}
\operatorname{map}(f,\operatorname{Failure}(e))=
\operatorname{Failure}(e).
\label{eq:map-failure}
\end{equation}
Both statements are \formalized{}.

\subsection{Horn support}

A finite Horn system is $H=(F_0,R,U)$, with base facts $F_0\subseteq U$ and rules whose premises and heads lie in $U$. Its immediate consequence operator is
\begin{equation}
T_H(S)=F_0\cup
\{\operatorname{head}(r)\mid r\in R,
\operatorname{prem}(r)\subseteq S\}.
\label{eq:horn-operator}
\end{equation}
Starting at the empty set gives the ascending chain
\begin{equation}
I_0=\varnothing,\qquad I_{n+1}=T_H(I_n),\qquad I_n\subseteq I_{n+1}.
\label{eq:horn-iteration}
\end{equation}
Finiteness yields stabilization by (|U|), and the closure
\begin{equation}
C_H=I_{|U|}
\label{eq:horn-closure}
\end{equation}
is a fixed point contained in every fixed point. These are \formalized{} results, consistent with the fixed-point basis of Horn semantics \citep{Horn1951,Tarski1955,CousotCousot1977}.

\subsection{Abstract argumentation}

An abstract argumentation framework is $AF=(A,D)$, with finite arguments $A$ and defeats $D\subseteq A\times A$ \citep{Dung1995}. The characteristic function is
\begin{equation}
F_{AF}(S)=\{a\in A\mid
\forall b\,((b,a)\in D\Rightarrow
\exists c\in S:(c,b)\in D)\}.
\label{eq:dung-characteristic}
\end{equation}
A set is conflict-free when none of its members defeats another. It is admissible when conflict-free and defending each member, complete when admissible and equal to (F_{AF}(S)), preferred when inclusion-maximal admissible, stable when it attacks every argument outside it, and grounded when it is the least fixed point. For finite (A), the implemented grounded extension is
\begin{equation}
G=F_{AF}^{|A|}(\varnothing),\qquad F_{AF}(G)=G.
\label{eq:grounded}
\end{equation}
These profile definitions and their finite results are \formalized{}. Which profile a legal institution should choose is not.

\subsection{Validation product}

The architecture distinguishes validation dimensions rather than collapsing them:
\begin{equation}
\mathcal V=
\mathcal V_{kernel}\times\mathcal V_{build}\times
\mathcal V_{runtime}\times\mathcal V_{legal}.
\label{eq:validation-product}
\end{equation}
Lean kernel acceptance, reproducible CI, executable refinement, and authorized legal judgment occupy different coordinates. There is no theorem of the form
$\mathcal V_{kernel}=1\Rightarrow\mathcal V_{legal}=1$. The absence of that implication is a scope rule, not a defect to be hidden.

Each coordinate has a different falsifier. A malformed proof term falsifies kernel acceptance. A missing import or stale generated artifact falsifies reproducible build evidence. A counterexample fixture falsifies a runtime refinement claim. A court, legislature, or authorized legal reviewer may reject a normative premise even when the preceding three coordinates are perfect. Treating these falsifiers separately avoids the category error in which success in one layer immunizes another.

The finite setting is equally important. Powerset enumeration and the $|A|$- or $|U|$-bounded iterations are executable reference semantics because their carriers are finite. They do not establish tractability for every production-sized instance. Complexity, indexing, incremental maintenance, and memory behavior are empirical engineering questions. The reference evaluator supplies an oracle against which optimizations can be tested, while Equation~\eqref{eq:incremental-correct} later states the equality an optimized implementation must meet.

Finally, the notation is intentionally extensional. Sets identify carriers by membership, and functions are judged by observable results. This keeps the mathematical interface small. Provenance, authority, and temporal annotations are nevertheless retained in the data types so that two extensionally equal payloads do not become interchangeable when their legal or evidential contexts differ.

\section{Typed Ontology, Evidence, and Execution}
\label{sec:ontology}

The operational ontology is a request-indexed typed graph. Let (G=(r_G,V_G,E_G)). An edge is well formed only when the graph request is itself well formed, the edge belongs to the graph, its request agrees with the graph, and all endpoints are present:
\begin{equation}
\operatorname{EdgeWF}(G,e)\iff
\operatorname{WF}(G.r)\land e\in E_G\land e.r=G.r
\land\operatorname{src}(e)\subseteq V_G
\land\operatorname{tgt}(e)\subseteq V_G.
\label{eq:edge-wf}
\end{equation}
The current formal node exposes an explicit kind tag. ULM03 does not prove a fully dependent payload refinement for every kind; any stronger ontology claim is \conjectured{}.

For state (s), an edge is enabled when its request agrees and its source nodes are active:
\begin{equation}
\operatorname{Enabled}(e,s)\iff e.r=s.r\land
\operatorname{src}(e)\subseteq s.active.
\label{eq:enabled}
\end{equation}
Application is additive:
\begin{equation}
\operatorname{applyEdge}(e,s).active=
s.active\cup\operatorname{tgt}(e).
\label{eq:apply-edge}
\end{equation}
A local transition exists when a well-formed enabled edge produces the target state. Request preservation for this relation is \formalized{}.

Edge kinds and claims determine required obligations. `typeSafety` is mandatory:
\begin{equation}
O_{req}(e)=\{\mathsf{typeSafety}\}\cup
\operatorname{flatMap}\bigl(O,
\operatorname{baseline}(e.kind)\mathbin{++}e.claims\bigr),
\quad O_{req}(e)\neq\varnothing.
\label{eq:required-obligations}
\end{equation}
The non-emptiness theorem is \formalized{}. It prevents an edge from bypassing all verification merely because its specialized claim list is empty.

A verifier advertises supported evidence kinds and is parameterized by a semantic guard (g):
\begin{equation}
\operatorname{VerifierSound}(g,v)\iff
\forall ev,\ v(ev)=\true\Rightarrow
ev.kind\in v.supported\land g(ev.subject).
\label{eq:verifier-sound}
\end{equation}
Satisfaction binds an obligation to evidence of the same subject and kind:
\begin{equation}
\operatorname{Sat}(v,s)\iff
s.o\in O_{req}(s.edge)\land
\exists ev,\ ev.subject=s\land ev.kind=s.o\land v(ev)=\true.
\label{eq:sat}
\end{equation}
From soundness and satisfaction, (g(s)) follows. This is \formalized{}. It is a contract theorem, not proof that an arbitrary Python or external verifier satisfies the premise.

Finite execution adds a completed-edge set:
\begin{equation}
\operatorname{runEnabled}(e,c)\iff e.r=c.r\land
\operatorname{src}(e)\subseteq c.active\land e\notin c.completed.
\label{eq:run-enabled}
\end{equation}
Application updates $active'=active\cup tgt(e)$, $completed'=completed\cup\{e\}$, and $phase'=phase+1$. The applied edge cannot immediately re-enable, and every finite run preserves request identity. Both are \formalized{}.

Finally, fact evidence retains the distinction between admissions and assumptions. For premise (p),
\begin{equation}
dep(p)=
\begin{cases}
\varnothing,&p=\mathsf{admitted}(a),\\
\{w.assumptionId\},&p=\mathsf{assumed}(w).
\end{cases}
\label{eq:premise-dependency}
\end{equation}
This \formalized{} dependency law prevents an assumed premise from being rendered as dependency-free fact. It does not determine whether an admission or assumption is legally justified.

The ontology therefore distinguishes typing from truth. A node can have the correct kind, request, and endpoints while carrying a false factual assertion. An edge can create exactly the mandated obligations while an external verifier remains unsound. The formal types make those risks visible and route them to the appropriate evidence gates; they do not erase the risks. This distinction is particularly important when language models propose nodes or edges. A syntactically valid proposal enters as candidate material and cannot become verified input until the advertised verifier contract is discharged.

Execution is additive only at the active-node and completed-edge level. The model does not infer that every enabled edge should run, that scheduling is fair, or that a production orchestrator terminates under arbitrary plugins. Those are separate operational properties. The finite-run theorem instead guarantees that whatever finite path is supplied cannot escape its request scope. That smaller theorem is compositional and reusable even when scheduling strategies change.

Provenance dependencies support negative auditing. If a conclusion depends on an assumption identifier, a renderer must retain that identifier or demonstrate an observation-preserving projection that intentionally summarizes it. Removing the identifier would not merely shorten an explanation; it would alter the observation being preserved. This creates a precise interface between proof objects and human-facing reports.

\subsection{Boundary countermodels and validation strategy}

Three small countermodels separate the ontology claims. First, consider an edge whose endpoints are present and whose transformation is otherwise sensible, but whose request differs from the graph request. Endpoint validation alone accepts the shape, while \operatorname{EdgeWF} rejects the composition. The failure shows why request identity is a semantic coordinate rather than decorative metadata. Second, consider an edge with no specialized claims. A checker that translates only those claims produces an empty worklist and may report vacuous success. The mandatory type-safety obligation prevents that path without pretending that type safety is the only required evidence. Third, consider a verifier that returns \true{} for every input. It can satisfy the syntax of a callback but cannot meet \operatorname{VerifierSound} for a nontrivial guard. Verifier success becomes meaningful only after the soundness premise is supplied.

These countermodels suggest a layered test design. Constructor tests should reject request and endpoint mismatch. Obligation tests should check that every edge kind produces the baseline obligation and that specialized claims add, rather than replace, evidence duties. Verifier tests should include a subject mismatch and an unsupported evidence kind. Execution tests should verify additive state change, completion marking, and request preservation along supplied finite paths. None of those executable checks proves the quantified Lean statements, but each can detect divergence between a concrete adapter and the formal interface.

The ontology remains intentionally incomplete in two respects. Node kinds are explicit, yet the payload type is not dependently indexed by every kind; a runtime can therefore still place semantically unsuitable content behind a syntactically correct tag. Scheduling is also left open: a finite run records chosen steps but does not prove fairness or maximal execution. Closing either gap would require new types and theorems, not stronger prose about the current declarations. Until then, payload validity and scheduler behavior belong in the open-obligation record.

\section{Finite Horn Closure and Incremental Support}
\label{sec:horn}

Horn support supplies the monotone base of the architecture. The immediate-consequence operator in Equation~\eqref{eq:horn-operator} is monotone because every premise inclusion that holds in (S) also holds in a superset (S'):
\begin{equation}
S\subseteq S'\Longrightarrow T_H(S)\subseteq T_H(S').
\label{eq:horn-monotone}
\end{equation}
This property is \formalized{}. It yields an ascending finite iteration rather than relying on an unbounded operational search. Since each $I_n\subseteq U$, either the chain adds a new element or it stabilizes. There can be at most $|U|$ strict additions, hence
\begin{equation}
I_{|U|}=I_{|U|+1}=T_H(I_{|U|}).
\label{eq:horn-stabilization}
\end{equation}

The fixed point is least. If $T_H(S)=S$, induction on $n$ gives $I_n\subseteq S$; therefore
\begin{equation}
C_H=I_{|U|}\subseteq S.
\label{eq:horn-least}
\end{equation}
The proof is machine checked in the finite setting. It justifies reading (C_H) as the supported fact carrier generated by the stated base and rules. It does not justify the base facts or rule encoding themselves.

A candidate (c) is well formed only if it is request-bound and supported by this closure:
\begin{equation}
\operatorname{CandidateWF}(H,c)\iff
c.r=H.r\land c.support\subseteq C_H.
\label{eq:candidate-wf}
\end{equation}
The formula is \formalized{}. It blocks a candidate from smuggling in a premise absent from the declared Horn system. It does not require that every fact in (C_H) be relevant to a candidate.

Incremental updates are restricted to additions. For $\Delta=(\Delta_F,\Delta_R)$,
\begin{equation}
H+\Delta=(F_0\cup\Delta_F,R\cup\Delta_R,U).
\label{eq:add-only-extension}
\end{equation}
The universe is fixed and the old facts and rules remain. Consequently,
\begin{equation}
T_H(S)\subseteq T_{H+\Delta}(S),
\qquad C_H\subseteq C_{H+\Delta}.
\label{eq:add-only-monotonicity}
\end{equation}
Both inclusions are \formalized{}. Retraction is excluded; a non-monotone change would require a different update theory and invalidation mechanism.

The repository specifies incremental implementation correctness extensionally:
\begin{equation}
\operatorname{IncrementalCorrect}(impl)\iff
\forall\Delta,\quad impl(\Delta)=
\operatorname{FullRecompute}(H+\Delta).
\label{eq:incremental-correct}
\end{equation}
This definition is \formalized{}, but it does not exhibit or verify a production worklist algorithm. A concrete implementation must provide a refinement witness or executable evidence. Confusing the contract with its implementation would reverse the direction of verification.

The separation between monotone support and later non-monotone acceptance is essential. New base facts or rules cannot erase previously supported facts under Equation~\eqref{eq:add-only-monotonicity}. Yet new attacks can change which arguments are accepted under an argumentation profile. The architecture therefore preserves the monotone carrier (C_H) while allowing the acceptance layer to expose disagreement rather than rewriting derivability.

This separation yields a useful diagnostic taxonomy. If an expected proposition is absent from $C_H$, the defect lies in the base facts, rule set, universe, or closure implementation. If the proposition is supported but no associated argument is accepted, the defect or disagreement lies in argument construction, attack generation, defeat policy, or extension selection. Debugging can therefore target a layer instead of treating every unfavorable result as a rule-engine failure.

The add-only restriction also defines a safe cache boundary. A cached old closure remains a subset of the new closure, but it is not necessarily equal to it. A runtime may reuse old support as a seed only if it recomputes the consequences of added facts and rules and produces the same carrier as full recomputation. Deletions, authority withdrawal, or temporal expiry do not satisfy this premise and must trigger invalidation rather than incremental addition.

From a legal perspective, monotonic support should not be confused with precedential entrenchment. A fact can remain derivable in the rule carrier while a later exception, attack, or procedural bar defeats the argument built from it. The two-layer design records both facts: the support did not disappear, and its decisional role changed.

\subsection{Proof anatomy and operational consequences}

The finite closure result depends on three distinct ingredients. Monotonicity orders the iterations, the finite universe bounds the number of strict additions, and the fixed-point argument identifies the stabilized carrier. Dropping any ingredient changes the claim. A monotone operator on an infinite universe need not stabilize after a known finite number of steps. A finite but non-monotone operator may oscillate. A stable implementation output need not be least unless it is related to the iteration from the empty carrier. The formal theorem packages these facts for the specified Horn operator; it is not a generic termination certificate for arbitrary rule engines.

This decomposition improves failure reports. If a runtime exceeds the expected bound, reviewers should ask whether the executable universe matches the formal finite carrier, whether rules introduce undeclared symbols, or whether the implementation is repeating work without changing the set. If the output stabilizes but disagrees with reference closure, the relevant issue is refinement rather than mathematical existence. If an add-only update appears to delete support, either the update violated the premise or the implementation did not realize the specified union. Each diagnosis points to a concrete witness rather than the generic statement that ``reasoning failed.''

The least-closure theorem is also deliberately neutral about legal relevance. It derives every fact licensed by the supplied Horn system, including facts that may never enter an accepted argument. A production implementation can optimize demand-driven evaluation, but the optimization must preserve the observable closure required at its interface. Performance evidence may justify the optimization; only equality with the reference semantics justifies substituting it in the assurance argument.

\section{Structured Arguments, Defeats, and Branch Queries}
\label{sec:aaf}

The argument layer turns supported facts into inspectable reasoning structures. For a structured argument (a), direct support is
\begin{equation}
p\prec_a q\iff\exists e\in E_a,
p\in\operatorname{prem}(e)\land\operatorname{concl}(e)=q.
\label{eq:direct-support}
\end{equation}
ULM08 requires this relation to be well founded, keeps all referenced nodes inside the carrier, binds the request, and preserves premise dependencies. These safeguards are \formalized{}. They prohibit cyclic or dangling proof objects within the defined normal form; they do not establish that a generated argument has captured every legally material reason.

Coverage is deliberately relative to a frozen expected carrier:
\begin{equation}
\operatorname{ArgumentCoverage}(Expected,Actual)
\iff Actual=Expected.
\label{eq:argument-coverage}
\end{equation}
The Boolean checker is sound and complete for this equality, and well-formedness transfers when the expected set has already been proved well formed. This is \formalized{}. The expected set is an input to the claim; the theorem is not a general completeness result for argument discovery.

Attacks must carry a nonempty witness and agree on request identity:
\begin{equation}
\operatorname{AttackWF}(x)\iff
x.witness\neq\epsilon\land x.attacker.r=x.target.r.
\label{eq:attack-wf}
\end{equation}
A policy $\pi$ resolves attacks into defeats:
\begin{equation}
D_{\pi}=\{(x.attacker,x.target)\mid
x\in Attacks\land\pi.succeeds(x)=\true\}.
\label{eq:resolved-defeat}
\end{equation}
Every member of $D_{\pi}$ has a same-request, well-formed source attack. This provenance property is \formalized{}. The function $\pi.succeeds$ is uninterpreted beyond its Boolean result. Legal correctness of priorities, burdens, exceptions, or value preferences requires an independently justified policy \citep{PrakkenSartor1997,ModgilPrakken2013,BenchCapon2003}.

The extension semantics are those introduced in Section~\ref{sec:preliminaries}. Conflict-free and admissible sets satisfy
\begin{align}
\operatorname{ConflictFree}(S)&\iff S\subseteq A\land
\forall a,b\in S,\ (a,b)\notin D,\label{eq:conflict-free}\\
\operatorname{Admissible}(S)&\iff
\operatorname{ConflictFree}(S)\land\forall a\in S,\operatorname{Defends}(S,a).
\label{eq:admissible}
\end{align}
Complete, preferred, stable, and grounded profiles are \formalized{}, including finite existence of preferred extensions and sound/complete enumeration relative to the defined carrier.

For profile (p), the extension family is explicit:
\begin{equation}
Ext_p(AF)=
\begin{cases}
\{G\},&p=\mathsf{grounded},\\
PreferredExt(AF),&p=\mathsf{preferred},\\
StableExt(AF),&p=\mathsf{stable},\\
CompleteExt(AF),&p=\mathsf{complete}.
\end{cases}
\label{eq:extension-family}
\end{equation}
Skeptical and credulous acceptance use different quantifiers:
\begin{align}
\operatorname{Common}(q)&\iff
\forall e\in Ext_p(AF),\operatorname{AcceptedIn}(e,q),
\label{eq:common}\\
\operatorname{Possible}(q)&\iff
\exists e\in Ext_p(AF),\operatorname{AcceptedIn}(e,q).
\label{eq:possible}
\end{align}
Corresponding refuted, possibly refuted, undecided, and inconsistent predicates are also defined. Query results retain branch identity, with
\begin{equation}
\operatorname{Composable}(x,y)\iff x.branch=y.branch.
\label{eq:branch-composable}
\end{equation}
All these statements are \formalized{}. Equation~\eqref{eq:branch-composable} prevents a report from selecting mutually favorable fragments from incompatible extensions. The model represents disagreement rather than manufacturing consensus.

Three counterexamples explain the boundary. First, a source-bound attack can still embody an incorrect priority rule; provenance is necessary but not sufficient for legal validity. Second, equality to a frozen expected carrier can pass even when the expected carrier omitted an argument; checker completeness is relative to its oracle. Third, a stable extension may fail to exist while complete or grounded semantics remain defined. A renderer that treats an empty stable-extension list as a negative legal answer would conflate semantic non-existence with claim rejection.

Branch identifiers also support explanation. A skeptical result may cite the intersection of all extensions, while a credulous result must identify at least one witness extension. Counterarguments should remain attached to the branch in which they operate. If two branches use the same textual argument identifier but different defeat relations, request and branch identity prevent their accidental coalescence.

Institutional policy enters at two explicit points: defeat resolution and profile selection. The formal layer can prove that the chosen Boolean policy is applied consistently and that the chosen profile satisfies its mathematical semantics. It cannot decide whether a burden of proof, priority rule, value ordering, or skeptical standard is legally required. Those choices must be supplied by an authorized source and retained in the evidence record.

\subsection{Branch-sensitive countermodels}

Suppose an argumentation framework has two incompatible preferred extensions. One accepts argument $a$ and rejects $b$; the other accepts $b$ and rejects $a$. A report that extracts the favorable conclusion from each extension can present $a$ and $b$ together even though no admissible branch supports that conjunction. This is not a failure of preferred semantics. It is a composition error introduced after evaluation. Recording the branch identifier in each query result makes the error observable, and the composability predicate refuses the join.

A second countermodel concerns quantifiers. If one extension accepts $a$ and another does not, credulous acceptance of $a$ is true while skeptical acceptance is false. Replacing the existential query with a universal one, or presenting either result without its profile, changes the proposition. The renderer must therefore preserve the query mode, semantic profile, and witness branch. Merely displaying the argument text is insufficient because the same text can have different acceptance status under different defeat relations or profiles.

A third countermodel separates provenance from legal validity. An attack can contain a nonempty witness, join arguments from the same request, and pass the formal well-formedness predicate while relying on an incorrect priority supplied by policy. The provenance theorem still holds: the resulting defeat came from a recorded attack and policy result. What fails is the external legal premise. The proper recovery is to replace or reauthorize the policy and reevaluate the branches, not to weaken the theorem that records where the defeat came from.

These examples determine an evaluation protocol. Tests should construct frameworks with multiple extensions, no stable extension, skeptical/credulous divergence, and policy-sensitive defeats. Each result should be serialized with request, profile, branch, and relevant witnesses, then compared with the finite reference semantics. The protocol evaluates executable fidelity to the formal model. It does not measure whether the arguments are persuasive to lawyers, whether the policy represents governing law, or whether the expected argument carrier is complete. Those questions require separate empirical and legal studies.

\section{Stratified Evaluation, Exact Calculation, and Conservative Assurance}
\label{sec:evaluator}

The evaluator does not map every input directly to a binary legal conclusion. It first preserves procedure. Applying a procedural cause changes the stage while leaving the normative marker invariant:
\begin{equation}
\operatorname{normativeMarker}(\operatorname{applyProcedureCause}(c,s))
=\operatorname{normativeMarker}(s).
\label{eq:procedure-preserves-marker}
\end{equation}
This is \formalized{}. It prevents procedural routing from silently rewriting substantive status.

The adjudication result has four constructors:
\begin{equation}
\operatorname{Adjudicate}\in
\{\mathsf{adjudicated},\mathsf{procedural},
\mathsf{pendingLegalJudgment},\mathsf{solverIncomplete}\}.
\label{eq:adjudication-four-way}
\end{equation}
An incomplete computation maps to `solverIncomplete`; a missing authority maps to `pendingLegalJudgment`; procedural material remains procedural; and substantive success or failure requires request-bound authority. These branch properties are \formalized{}. They refuse the common but invalid inference from “the solver returned no positive result” to “the legal claim fails.”

Domain composition selects nonempty candidate subsets under an explicit policy:
\begin{equation}
Choices(B,\pi)=\{S\subseteq B.candidates\mid
S\neq\varnothing\land\pi(S)=\true\}.
\label{eq:composition-choices}
\end{equation}
Branch, policy identifier, and selected subset are part of the resulting identity. This \formalized{} construction makes the choice auditable but does not prove that $\pi$ is the correct legal policy.

Exact quantities use a dimension index. Their denotation obeys
\begin{align}
\llbracket lit(q)\rrbracket&=q,\label{eq:eval-lit}\\
\llbracket x+y\rrbracket&=\llbracket x\rrbracket+\llbracket y\rrbracket,
\label{eq:eval-add}\\
\llbracket x-y\rrbracket&=\llbracket x\rrbracket-\llbracket y\rrbracket,
\label{eq:eval-sub}\\
\llbracket qx\rrbracket&=q\llbracket x\rrbracket.
\label{eq:eval-scale}
\end{align}
Type indices reject direct addition across incompatible currencies, periods, or other dimensions. The executable evaluator equals the recursive denotation. This is \formalized{} arithmetic correctness, not proof of the statutory interpretation that selected the expression.

Coverage is complete exactly when no obligations remain:
\begin{equation}
\operatorname{Complete}(C)\iff C.openObligations=\varnothing.
\label{eq:coverage-complete}
\end{equation}
Composition cannot erase obligations:
\begin{equation}
open(C_1\sqcup C_2)=open(C_1)\cup open(C_2).
\label{eq:coverage-union}
\end{equation}
Both are \formalized{}.

Trust is an ordinal five-vector rather than a calibrated probability:
\begin{equation}
\tau=(\tau_{source},\tau_{text},\tau_{fact},\tau_{proof},\tau_{authority})
\in\{0,1,2\}^{5}.
\label{eq:trust-vector}
\end{equation}
The coordinates record source, text, fact, proof, and authority assurance. Their meet is coordinatewise minimum:
\begin{equation}
(\tau\wedge\sigma)_i=\min(\tau_i,\sigma_i),
\qquad\tau\wedge\sigma\preceq\tau,\sigma.
\label{eq:trust-meet}
\end{equation}
The meet laws are \formalized{}. Assurance records combine only under the same scope, select the weaker implementation and run statuses, and union obligations, legal inputs, trusted-computing-base references, and notices. Aggregation may preserve or reduce assurance; it cannot improve it by discarding the weaker component.

The evaluator is stratified because each stage has a distinct blocking condition. A solver may be complete while authority is missing; an authority may be valid while evidence obligations remain open; an exact arithmetic expression may evaluate while its legal basis is disputed. A single Boolean would erase these combinations. The product record retains them so a downstream user can distinguish technical remediation from legal review.

Exact dimensional types are especially useful for negative guarantees. They reject nonsensical additions before evaluation, but they cannot infer exchange rates, statutory interest periods, rounding regimes, or damages rules. Those values and policies must enter as typed, sourced inputs. Once supplied, denotational equality shows that the evaluator implements the stated expression exactly.

The five trust coordinates should not be averaged. An arithmetic mean would allow high proof assurance to offset absent authority or weak factual and textual support. Coordinatewise meet implements a non-compensation principle: every dimension remains separately visible, and the weakest input bounds the combined result in that dimension. Whether the ordinal values $0,1,2$ should have a different institutional interpretation is a governance choice; their algebraic meet is the proved part.

\subsection{Failure localization as an evaluation method}

The four-way adjudication result supplies a diagnostic protocol. A \mathsf{solverIncomplete} value asks for the listed computational obligations, not a new legal conclusion. A \mathsf{pendingLegalJudgment} value asks for qualified authority or a missing finding, not another solver run. A procedural disposition identifies a routing state whose normative marker has not been rewritten. Only an adjudicated status represents the modeled substantive branch, and even then its validity is conditional on the supplied facts, norms, profile, and authority object. This stratification prevents remediation at one layer from being misreported as resolution at another.

Exact arithmetic has a similar boundary. Dimension indices can reject addition of incompatible quantities and the evaluator can be shown equal to its denotation. Neither result determines the economically or legally correct operands. A statutory-interest calculation, for example, still needs an authorized principal, rate, period, day-count convention, and rounding rule. If one input changes, the request or evidence digest should change as well. The mathematical evaluator can then reproduce the declared expression while the legal provenance of each operand remains separately reviewable.

Trust composition should be tested coordinate by coordinate. A useful negative fixture pairs an assurance record with strong source and proof coordinates but weak fact authority with another record having the inverse pattern. The meet must retain both weaknesses rather than averaging them into a moderate score. Coverage fixtures should likewise show that an obligation appearing in either input remains open after composition. These are executable tests of the record operations; they do not calibrate what levels zero, one, and two mean in a particular institution.

Temporal and procedural properties remain outside the five fields even when a surrounding assurance envelope records them elsewhere. Calling a coordinate ``temporal'' or ``procedural'' would misstate the formal \operatorname{TrustVector} type. A future extension could introduce additional coordinates, but it would need new meet definitions, ordering laws, migration rules, and evidence about how old records map into the expanded carrier. The current result is limited to the five named fields above.

\section{Temporal Applicability}
\label{sec:temporal}

Legal authority is versioned, and observations must be evaluated relative to an as-of time. Closed intervals intersect according to
\begin{equation}
[a,b]\cap[c,d]=
\begin{cases}
[\max(a,c),\min(b,d)],&\max(a,c)\le\min(b,d),\\
\varnothing,&\text{otherwise}.
\end{cases}
\label{eq:interval-intersection}
\end{equation}
This arithmetic is \formalized{} for the repository's interval type.

For a version (v), applicability at time (t) is
\begin{equation}
\operatorname{effectiveAt}(v,t)\iff
v.from\le t\land(v.to=\bot\lor t\le v.to).
\label{eq:effective-at}
\end{equation}
Observation respects the query cut-off only when
\begin{equation}
\operatorname{observationAllowed}(t_{obs},t_{asof})
\iff t_{obs}\le t_{asof}.
\label{eq:asof}
\end{equation}
Future observations, retracted versions, and superseded sources are rejected. These predicates and rejection properties are \formalized{}.

The Kripke component uses a reachability relation (R^+). An always modality is read as
\begin{equation}
K,i\models G\varphi\iff
\varphi(i)\land\forall j\,(iR^+j\Rightarrow\varphi(j)).
\label{eq:always}
\end{equation}
When every reachable world satisfies the required temporal ordering, the global guard follows. This theorem is \formalized{} and belongs to the modal and temporal tradition \citep{Kripke1963,Pnueli1977,ClarkeEtAl1986}. Its premise already quantifies over the worlds. It is not a transition-invariant theorem that derives future safety from an initial state alone.

The legal interpretation is limited. A timestamp and version interval can prevent use of information outside the declared as-of boundary. They cannot establish publication authenticity, jurisdictional applicability, or correct interpretation. Those properties must be carried as evidence obligations and authority decisions.

Temporal failures should be classified rather than flattened. A future observation violates the query cut-off; an expired version falls outside its validity interval; a retracted source is disabled by status; and a superseded source may remain historically true but no longer govern. Each case demands a different explanation and recovery action. Moving the as-of date, selecting an earlier version, or obtaining a new authority record are not interchangeable operations.

The model also distinguishes event time from processing time. A document retrieved today may describe an event that occurred before the as-of date, but its later publication may make it unavailable to the historical decision-maker. A complete epistemic model would carry both coordinates and a rule for admissible knowledge. The current ULM predicates formalize the stated observation time and version checks; richer epistemic availability remains \conjectured{}.

For reproducibility, a temporal query must bind the authority version, as-of instant, observation cut-off, and timezone or calendar convention. Without those fields, repeating the same textual query can select different legal materials while appearing identical. Request-bound temporal metadata turns that drift into an observable input change.

\subsection{Temporal countermodels and recovery}

Consider first a source that was effective at the event date but retracted before the query date. Interval membership alone may accept it, while the status predicate rejects it. Consider next a later source that accurately summarizes an earlier event but was unavailable at the requested historical cut-off. Event time alone may accept the content, while observation time rejects its use in reconstructing the earlier knowledge state. Finally, consider two formally identical rules issued by different authorities or for different jurisdictions. Matching text and dates do not establish applicability because scope and authority remain independent coordinates.

These countermodels explain why time cannot be reduced to a single timestamp. Version validity, observation availability, event occurrence, and processing time answer different questions. The current formal predicates cover the declared version and observation checks. A claim about what a decision maker actually knew would require a richer epistemic model and evidence about publication, access, and notice. Such a claim is \conjectured{} here and must not be inferred from the always-modality theorem.

Testing should vary one temporal field at a time. Boundary fixtures include an observation exactly at the as-of instant, one just after it, an open-ended version, a closed version at its terminal date, an expired version, a retracted version, and a superseded source retained only for historical explanation. Expected outcomes should record both rejection reason and recovery action. Changing an as-of date is a new query, not a repair of the old one; selecting a different version requires new provenance; restoring a retracted source requires authority evidence rather than clock arithmetic.

The modality has a further proof boundary. Its premise already asserts the guarded property for every reachable world. It does not prove that an arbitrary transition relation preserves the guard, that new worlds cannot be added, or that the initial state alone determines future safety. A production system claiming temporal invariance must separately establish transition preservation and relate its evolving state graph to the modeled reachability relation.

\section{A Compositional Rosetta Stone}
\label{sec:rosetta}

The architecture links representations through preservation laws rather than asserting a category-theoretic equivalence. For stages (A,B,C), observation-preserving maps satisfy
\begin{equation}
P(f)\land P(g)\Longrightarrow P(g\circ f),
\label{eq:preservation-composition}
\end{equation}
where (P(f)) abbreviates Equation~\eqref{eq:observation-preservation}. This is \formalized{}.

The same-request condition is repeated at every composition boundary. For branch results,
\begin{equation}
x\bowtie y\iff x.branch=y.branch.
\label{eq:branch-join}
\end{equation}
For assurance records,
\begin{equation}
A\oplus B\text{ is defined only if }scope(A)=scope(B).
\label{eq:scope-join}
\end{equation}
For domain choices,
\begin{equation}
id(choice)=\langle request,branch,policyId,selectedSet\rangle.
\label{eq:choice-identity}
\end{equation}
These constraints are \formalized{} in their respective types. Together they provide a typed correspondence among request normalization, support, argumentation, procedure, and assurance.

It is useful to express the intended pipeline as a partial composite
\begin{equation}
\mathcal L=
Release\circ Assure\circ Adjudicate\circ Query\circ
EvaluateAF\circ CompileArgs\circ CloseHorn\circ Normalize.
\label{eq:pipeline-composite}
\end{equation}
This equation is \derived{} notation. The Lean corpus proves properties of the named stages and selected compositions; it does not define a category, functor, natural transformation, or equivalence covering the full pipeline. Any such categorical account is \conjectured{} until objects, morphisms, identities, composition laws, and preservation theorems are explicitly encoded.

The Rosetta-stone value is practical rather than ornamental. A runtime may use JSON, a theorem may use an inductive type, and a report may use prose. Observation-preservation and shared identity fields state what must survive each representation change. The approach avoids an unjustified claim that syntactically different artifacts are semantically identical in every respect.

There are two distinct composition failures. A structural failure occurs when scopes, branches, or request identifiers differ, so the operation is undefined. An evidential failure occurs when structures align but an observation-preservation or verifier obligation remains open. The first should stop the composition immediately; the second may yield a partial result with explicit obligations. Collapsing both into a generic error would lose useful recovery information.

The partial composite in Equation~\eqref{eq:pipeline-composite} also reveals where alternative implementations can enter. A different Horn engine, argument generator, extension enumerator, or renderer may replace a stage if it proves or tests the same observable contract. Internal algorithms need not be identical. What must remain stable is the request-bound behavior declared at the interface.

This is why the architecture prefers small preservation theorems to one grand equivalence claim. A global equivalence would require a semantics for every representation and every loss of detail. Local observations let the specification name only the information that a transition promises to preserve, while open obligations expose everything not yet justified.

\subsection{Interface-by-interface verification}

Each arrow in Equation~\eqref{eq:pipeline-composite} needs its own witness. Normalization should preserve the selected observation of a request. Horn closure should preserve the request and expose its finite support carrier. Argument compilation should bind premises, dependencies, and conclusions to that carrier. Abstract evaluation should retain the defeat policy and semantic profile. Querying should attach the branch and quantifier mode. Adjudication should distinguish computational incompleteness, procedural disposition, absent authority, and modeled substantive status. Assurance should join only compatible scopes and retain weaker coordinates. Release should verify the subject-bound evidence tuple. No witness at a later arrow discharges a missing witness earlier in the sequence.

This approach permits implementation diversity without semantic free choice. A JSON adapter and a Lean constructor may use different layouts, yet an observation-preservation test can compare their request, scope, and obligation projections. An optimized extension enumerator may use pruning rather than powerset enumeration, yet it must agree with the reference extension family for the tested finite carrier. A report generator may shorten internal identifiers, yet it must retain a reversible link when those identifiers are part of the promised observation. The interface specifies what equivalence means for that transition.

Lossy transformations are not automatically forbidden. They are acceptable when the lost information is outside the declared observation and when the decision to omit it is itself authorized. A public explanation may redact sensitive evidence while retaining a sealed provenance reference. That operation is not equality of complete artifacts; it is preservation of a deliberately narrower observation plus an outstanding access or review obligation. The formal composition theorem can apply only after that observation has been defined.

The partiality of the composite is therefore informative. Scope mismatch should produce no composition. Missing evidence may produce an explicit partial outcome. A runtime exception produces failure. A legally contested policy may leave technically complete results pending judgment. These outcomes have different recovery paths, and the Rosetta stone is successful only if translation preserves those distinctions rather than normalizing all of them to a Boolean error or success flag.

\section{Conditional Convergence and Weighted Distance}
\label{sec:banach}

An empirical attachment retains normative solutions unchanged while adding observational material. A weighted deviation score is
\begin{equation}
D(w,x)=\sum_{i=1}^{n}w_i x_i.
\label{eq:deviation-score}
\end{equation}
Both read-only attachment and the score definition are \formalized{}. The score is not a calibrated probability or legal burden of proof.

For positive weights (w_i), define the weighted sup distance on finite real vectors:
\begin{equation}
d_w(x,y)=\max_i\frac{|x_i-y_i|}{w_i},\qquad w_i>0.
\label{eq:weighted-sup}
\end{equation}
Non-negativity, symmetry, the triangle inequality, and point separation are \formalized{}. Despite one historical declaration name, the source does not establish a `CompleteSpace` instance for this distance. Completeness must be provided separately before invoking Banach.

Suppose (T) has coordinatewise Lipschitz bounds
\begin{equation}
|T(x)_i-T(y)_i|\le\sum_j L_{ij}|x_j-y_j|
\label{eq:coordinate-lipschitz}
\end{equation}
and the weighted row condition
\begin{equation}
\sum_jL_{ij}w_j\le q w_i.
\label{eq:weighted-row}
\end{equation}
Then
\begin{equation}
d_w(Tx,Ty)\le qd_w(x,y).
\label{eq:weighted-contraction}
\end{equation}
The implication is \formalized{}. If in addition $0\le q<1$, the space is nonempty and complete, and $T$ is a contraction, Banach's theorem \citep{Banach1922} gives a unique fixed point $x^*$ and
\begin{equation}
d(T^n x,x^*)\le\frac{q^n}{1-q}d(x,Tx).
\label{eq:banach-error}
\end{equation}
This general theorem is \formalized{}.

No source theorem shows that the full legal evaluator satisfies these premises. Mapping a legal graph into $\mathbb R^n$, choosing weights, proving completeness, and establishing $q<1$ are additional burdens. Convergence of an iterative score is also not convergence to legal truth. Those applications are \conjectured{}.

The distinction between a metric and a ranking score is decisive. A symmetric nonnegative score may still fail identity of indiscernibles or the triangle inequality. If distinct graphs receive distance zero, a fixed-point argument on the resulting representation may converge only up to an equivalence class. A quotient construction could repair point separation, but then the legal meaning of the quotient must be justified.

Weights are likewise normative or empirical inputs. Positive weights make Equation~\eqref{eq:weighted-sup} well-defined; they do not show that one legal feature deserves twice the influence of another. Sensitivity analysis should therefore vary weights and report whether conclusions are stable. A formal contraction proof under fixed weights cannot replace that model-selection inquiry.

An executable convergence claim needs four independent witnesses: the state-space embedding, metric-space completeness, a verified Lipschitz bound, and correspondence between the implemented iteration and the formal map. Missing any witness leaves the Banach result inapplicable. The release record should state which premise failed instead of reporting merely that convergence was not certified.

\subsection{Counterexamples to premature convergence claims}

Several near misses clarify the premises. If a proposed distance maps two legally distinct graphs to the same feature vector, point separation fails even when symmetry and the triangle inequality hold on vectors. The iteration may converge in feature space while alternating between distinct legal structures hidden by the embedding. If some weight is zero, Equation~\eqref{eq:weighted-sup} is undefined; if a weight is negative, it no longer supports the intended metric. If the weighted row bound holds only on sampled pairs, the observation is evidence for a conjectured Lipschitz property rather than a quantified proof.

The constant matters independently. A map with $q=1$ can be nonexpansive without having a unique fixed point, and a map with a local $q<1$ may leave the region on which the bound was measured. Completeness also cannot be borrowed from an ambient space when the actual state subset is not closed. These mathematical distinctions are familiar, but they are easy to lose when a numerical process appears stable in a finite trace. Empirical stabilization is neither the contraction premise nor the uniqueness conclusion.

A defensible evaluation separates four tasks. First, define an embedding and test whether known legally material distinctions survive it. Second, verify the metric laws for the actual carrier or state the quotient being used. Third, establish the row or Lipschitz bound with a proof or a clearly labeled empirical estimate. Fourth, compare executable iterations with the formal map and report residual error under the same distance. Only after all four tasks succeed may Banach's theorem be applied, and even then the fixed point is a fixed point of the modeled transformation rather than a legally correct outcome.

Failure recovery follows the missing premise. Colliding embeddings require additional features or an explicit quotient. Invalid weights require redefining the distance. An unproved global bound requires a restricted domain or a weaker convergence claim. Runtime disagreement requires refinement work. None of these repairs can be achieved by renaming a similarity score as a metric or by interpreting a stable output as legal truth.

\section{Differential Privacy and the Privilege Underdetermination Result}
\label{sec:dp}

Differential privacy is a quantitative property of a randomized mechanism, not a label attached directly to a legal category. For an adjacency relation $\sim$, mechanism $M$ is $(\varepsilon,\delta)$-differentially private when
\begin{equation}
\Pr[M(D)\in S]\le e^{\varepsilon}\Pr[M(D')\in S]+\delta
\label{eq:dp}
\end{equation}
for all $D\sim D'$ and measurable $S$ \citep{DworkEtAl2006,DworkRoth2014}. Sensitivity for query $f$ is
\begin{equation}
\Delta f=\sup_{D\sim D'}\|f(D)-f(D')\|.
\label{eq:sensitivity}
\end{equation}
The scale of a Laplace mechanism depends on both $\Delta f$ and a chosen $\varepsilon$, not on a privilege classification alone.

Let (P) be a legal privilege label. The rejected shortcut is a universal function
\begin{equation}
\varepsilon=\phi(P).
\label{eq:privilege-epsilon}
\end{equation}
The same privilege category can be paired with different adjacency relations, threat models, queries, utility constraints, composition horizons, and institutional risk tolerances. Therefore (P) does not determine a unique privacy budget. The repository records this as a refuted theory claim supported by two-model counterexample evidence; it is not represented as a Lean differential-privacy theorem. We label the underdetermination result \derived{}, not \formalized{}.

For repeated mechanisms, a basic composition bound has the form
\begin{equation}
(\varepsilon_1,\delta_1)\oplus(\varepsilon_2,\delta_2)
=(\varepsilon_1+\varepsilon_2,\delta_1+\delta_2).
\label{eq:basic-composition}
\end{equation}
More refined accounting and smooth sensitivity \citep{NissimEtAl2007} do not change the conceptual point: legal authorization constrains permissible processing, while privacy parameters quantify properties of a specified mechanism. Neither coordinate replaces the other.

The ULM trust vector and authority receipts can carry privacy-related assurance without fabricating a budget. If $\tau_p$ records procedural or privacy review, conservative composition uses
\begin{equation}
(\tau\wedge\sigma)_p=\min(\tau_p,\sigma_p).
\label{eq:privacy-trust-meet}
\end{equation}
This is \formalized{} as part of the general trust meet. It is ordinal metadata, not a $\varepsilon$ estimate.

A future extension must define adjacency, mechanism semantics, budget accounting, and authority rules, then prove a connection between executable code and Equation~\eqref{eq:dp}. Until then, any automatic privilege-to-budget mapping is \conjectured{} and must not appear as a formal release claim.

The two-model counterexample is structural. Hold the privilege label fixed while changing either the adjacency relation or the query sensitivity. One admissible institution may choose a narrow neighboring-dataset relation and a moderate utility target; another may choose a broader relation and stricter disclosure risk. Both respect the same legal category yet require different mechanism parameters. Therefore uniqueness cannot follow from the label alone.

Authorization and privacy proof should be represented as a product rather than a scalar. The first component states that an actor may perform a purpose-bound operation; the second states the quantitative guarantee of the mechanism under declared adjacency and composition assumptions. Authorization without a mechanism proof leaves technical privacy open. A mechanism proof without authorization leaves legal permissibility open.

Failure recovery follows the missing component. An undefined adjacency relation requires a data-governance decision; an unbounded sensitivity requires redesign or clipping; exhausted composition budget requires stopping or new authorization; a receipt mismatch requires renewed human review. None of these defects can be repaired by relabeling the data as privileged.

\subsection{Two-model construction and evaluation consequences}

The underdetermination argument can be made without choosing numerical budgets. Fix one privilege label $P$ and construct Model A with record-level adjacency, a bounded count query, a single release, and one stated utility constraint. Construct Model B with user-level adjacency, a vector query of different sensitivity, repeated releases, and a different threat model. Both models classify the underlying material by $P$, yet the feasible mechanism families and accounting obligations differ. If a function $\phi(P)$ returned one uniquely justified budget, it would have to ignore those differences. The contradiction concerns sufficiency of the legal label, not the existence of an institutionally selected budget after the missing parameters are supplied.

A second construction keeps adjacency fixed and changes the authorized purpose. The same mechanism may be permitted for an internal audit but not for public release, or may require different composition horizons. Differential privacy alone does not answer that authorization question. Conversely, authorization to process confidential material does not establish Equation~\eqref{eq:dp}. The two properties can coexist, fail independently, and require different evidence. The appropriate record therefore carries both rather than converting one into the other.

Evaluation should report adjacency, query, sensitivity method, mechanism, $\varepsilon$, $\delta$, composition accountant, number of releases, clipping or truncation rules, utility criterion, and approving authority. Omitting any field limits what can be reproduced. A mechanism test can inspect distributional behavior on selected neighboring datasets, but finite sampling does not prove the universal probability inequality. A formal proof can establish the inequality for a model while leaving the deployed randomness, serialization, or accounting implementation unverified. Runtime refinement and statistical tests remain separate evidence layers.

Recovery must also preserve prior releases. Lowering a future budget does not undo privacy loss already composed, and changing adjacency changes the statement being protected rather than repairing the original statement. If an accountant is wrong, the certificate for affected releases should be blocked or revoked and recomputed from retained events. If authorization expires, technically valid privacy parameters do not extend it. These cases illustrate the paper's broader thesis: assurance coordinates compose conservatively and cannot substitute for one another.

\section{Authority and Taint Non-Interference}
\label{sec:noninterference}

Authorization is modeled as a ranked capability rather than a vote count. For level (l) and action kind (k),
\begin{equation}
\operatorname{canIssue}(l,k)\iff
rank(l)\ge rank(required(k)).
\label{eq:can-issue}
\end{equation}
A valid one-level receipt satisfies
\begin{equation}
\operatorname{receiptValid}(r)\iff
rank(r.to)=rank(r.from)+1.
\label{eq:receipt-valid}
\end{equation}
Same-level consensus cannot create a higher rank:
\begin{equation}
consensusRank([l]^n)\le rank(l).
\label{eq:consensus-rank}
\end{equation}
These results are \formalized{}. They directly address multi-agent systems: adding agents at the same authority level may diversify proposals but cannot authorize a transition reserved for a human or institution.

A human-research receipt is bound to a task and digest:
\begin{equation}
bind(r,t,d)\iff r.task=t\land r.digest=d.
\label{eq:receipt-bind}
\end{equation}
It is temporally valid only when
\begin{equation}
valid(r,now)\iff r.issued\le now\le r.expiry\land\neg r.revoked.
\label{eq:receipt-time}
\end{equation}
Cross-task, cross-input, expired, and revoked receipts fail. These are \formalized{} properties. A human-research receipt records an action and, when the relevant predicates hold, proves binding to the stated task and digest together with temporal validity and non-revocation. It does not by itself prove that the reviewer was legally authorized or competent, or that the substantive action was correct.

Information quality uses a two-point taint lattice:
\begin{equation}
\mathsf{clean}\sqcup\mathsf{clean}=\mathsf{clean},
\qquad x\sqcup\mathsf{tainted}=\mathsf{tainted}.
\label{eq:taint-join}
\end{equation}
For the formal stage constructor, a tainted head input propagates through the input fold:
\begin{equation}
taint(x)=\mathsf{tainted}
\Longrightarrow
taint(stageOutput(x::xs,c))=\mathsf{tainted}.
\label{eq:taint-propagation}
\end{equation}
The join laws, this \operatorname{stageOutput} result, and the majority non-washing results are \formalized{}. An arbitrary external combination function satisfies the same implication only if it refines this taint fold or supplies its own preservation proof. Duplicating an unsupported answer does not produce provenance:
\begin{equation}
\bigsqcup_{i=1}^{n}\mathsf{tainted}=\mathsf{tainted}.
\label{eq:no-taint-washing}
\end{equation}

Authority and taint are separate coordinates. A high-authority actor can authorize review of tainted material without making the material clean; clean machine evidence can exist without authority to issue a legal result. Conservative assurance carries both dimensions and prevents either from substituting for the other.

Multi-agent orchestration therefore needs independence claims that are stronger than agent count. If agents share a model, prompt, retrieval cache, or contaminated source, their errors may be correlated. The taint join remains conservative without estimating that correlation. A future probabilistic model could quantify dependence, but it must not weaken the deterministic rule that known taint survives aggregation.

Receipt chains also require continuity. A downstream action should identify the immediately preceding authority level, the same task and digest, and a valid time interval. Skipping a level or replaying a receipt from another task fails structurally. This converts an otherwise informal approval history into an inspectable path while leaving the substance of each human decision outside the theorem.

Recovery is explicit. Tainted evidence must be replaced or independently verified; an expired receipt must be reissued; a digest mismatch requires review of the actual new material; and insufficient authority requires escalation to the designated level. Adding more low-authority agents is not a recovery action. These rules preserve the distinction between search capacity, evidence quality, and decisional power.

\subsection{What the non-interference claims exclude}

The authority results constrain rank transitions, not institutional legitimacy. The \operatorname{AuthorityReceipt} structure records a subject, case scope, issuer, and adjacent levels, while its \operatorname{receiptValid} predicate checks the one-rank succession equation. The presence of identity fields in a structure does not mean that this predicate authenticates them. Likewise, the human-research receipt binds task and digest and checks a time window and revocation flag, but external authorization and reviewer competence remain premises for a legal process. A release record should name the predicate actually checked instead of summarizing every field as ``validated.''

The taint theorem is also intensional to the defined fold. It proves that \operatorname{stageOutput} computes the join of its input taints and that a tainted leading input makes the result tainted. List-level repetition and majority examples reuse that operation. The theorem does not quantify over every possible merger, voting scheme, language-model prompt, or database query. A concrete aggregator belongs inside the proved boundary only after its output label is related to \operatorname{taintOfInputs}.

These limits support precise negative tests. A rank-skipping authority receipt should fail succession. A receipt replayed with another task or digest should fail binding. An expired or revoked human receipt should fail current validity. A stage with one tainted input should remain tainted, including when clean inputs or repeated copies are added. An external aggregator that drops labels should fail refinement even if its textual conclusion matches a clean baseline. Passing these fixtures supplies bounded runtime evidence; it does not broaden the quantified Lean theorem.

The model intentionally does not define a procedure by which tainted material becomes clean. Independent verification may create a new clean evidence object, but relabeling the original object would erase provenance. The new object should record the verification method and its relationship to the source. Similarly, higher authority may approve use of uncertain material without changing its epistemic quality. Keeping those operations distinct prevents authorization from laundering evidence and prevents clean evidence from manufacturing legal power.

\section{Evidence-Calibrated Release}
\label{sec:evidence}

The theorem corpus and the release certificate answer different questions. A theorem is checked by the Lean kernel under its imports. A release certificate asserts that a named source subject has passed a specified collection of build and evidence gates. The evidence tuple is
\begin{equation}
E=\langle sha,tree,run,modules,cleanBuild,axioms,claims,
mutation,receipts,certificate,verdict\rangle.
\label{eq:evidence-tuple}
\end{equation}
The components must refer to one subject. A green job icon is not a substitute for inspecting their contents.

The formal inventory contains 145 declarations in the comprehensive audit and 27 declarations in the core-composition audit. The accepted axiom output is restricted to the standard dependencies
\begin{equation}
A_{allowed}=\{\mathsf{propext},\mathsf{Classical.choice},\mathsf{Quot.sound}\}.
\label{eq:allowed-axioms}
\end{equation}
No custom axiom is part of the audited proof boundary. These are commit-bound evidence facts; they must be regenerated after source changes.

Mutation testing supplies executable negative evidence. Let $b$ be a baseline fixture and $m_1,\ldots,m_k$ be property-breaking mutations. A genuine light gate requires
\begin{equation}
Gate(b,m_1,\ldots,m_k)=PASS
\iff Accept(b)\land\bigwedge_{i=1}^{k}\neg Accept(m_i).
\label{eq:mutation-gate}
\end{equation}
This criterion is \derived{} from the gate specification. It ensures that the report is generated by executions rather than fabricated as a static file. It does not establish completeness over all faults.

Cross-repository refinement receipts connect abstract verifier contracts to executable fixtures. For receipt (r), the required binding is
\begin{equation}
Binding(r)=\langle producerRepo,producerSha,consumerRepo,
consumerSha,task,digest,issued,expiry,authority\rangle.
\label{eq:cross-repo-binding}
\end{equation}
The ULM level formalizes narrower receipt components: one-rank authority succession and separate task, digest, time-window, and revocation predicates. Producer and consumer repository names and source SHAs in Equation~\eqref{eq:cross-repo-binding} are runtime-schema fields checked by the cross-repository pipeline, not fields proved by one ULM receipt theorem. A concrete receipt is therefore bounded runtime evidence. Equality between selected executable outputs and expected formal fixtures is not a universal proof of runtime equivalence.

Certificate generation is fail-closed. Missing required evidence produces a blocked status, generation or verification failure maps to job failure, and the independent verifier must agree with the certificate. Abstractly,
\begin{equation}
Verified(E)\iff SameSubject(E)\land CompleteEvidence(E)
\land CertificatePass(E)\land VerifierAgrees(E).
\label{eq:verified-release}
\end{equation}
This is a release-policy specification, not a Lean theorem. The machine-readable JSON artifacts, not this narrative, decide whether a particular subject satisfies it.

The release boundary also records what the formal development does not establish. Probability kernels, Bayesian calibration, generic analogy strength, a graph-similarity metric, human explanation quality, differential-privacy parameters, and substantive liability allocation are absent from the proved core. They may use the architecture's obligations and assurance types, but they require new semantics and evidence before receiving a \formalized{} label.

\subsection{Research integrity and reproducibility}

No private legal matter or customer dataset was used. Bibliographic entries were checked against DOI resolvers, publisher records, official journal sources, or EUR-Lex. The formal source remains primary for theorem claims. The recommended reproduction record is the exact commit, tree, workflow run, artifact names, certificate subject and status, missing-evidence array, and independent-verifier verdict.

The mutation gate and refinement receipts address different failure modes. Mutation asks whether a checker rejects deliberately invalid variants. Refinement asks whether an executable producer agrees with specified fixtures and binds its output to a source version. A system can pass one and fail the other. The certificate must therefore carry both results rather than replacing them with a single aggregate percentage.

Independent verification should recompute relations among artifacts, not merely trust their filenames. It checks subject identity, required fields, evidence status, and consistency between certificate and verifier verdict. When evidence is missing, the correct outcome is blocked rather than a partial pass disguised as success. This is fail-closed at the release layer in the same sense that `Outcome.Partial` and `solverIncomplete` remain explicit inside the model.

Evidence retention must also preserve negative results. Failed mutations, counterexamples to graph metrics, and the refuted privilege-to-privacy mapping are part of the scientific record. Deleting them would make later positive claims impossible to audit. Release documentation should state whether a limitation is unresolved, accepted as outside scope, or closed by a new theorem or fixture.

\subsection{Interpreting finite release evidence}

The evidence tuple should be read conjunctively and subject-specifically. A theorem audit establishes which declarations were checked under a particular source tree. A clean build establishes that the recorded environment compiled that tree. A mutation report establishes that enumerated property-breaking variants were rejected by the exercised checker. A refinement receipt establishes that selected executable fixtures produced the expected observations under named producer and consumer versions. The certificate and independent verifier then check consistency of those records. None of the components subsumes another, and none transfers automatically to a later commit.

Counts require the same discipline. A report that all designed mutations were rejected states coverage of the listed mutation catalogue, not coverage of all implementation defects. A report that three refinement fixtures passed states three witnessed executions, not three independent proofs or universal behavioral equivalence. The denominator must remain visible because ``all passed'' without the fixture inventory is not reproducible. Similarly, the 145- and 27-declaration audits are inventory facts for the reported subject; they are not permanent measures of project size or theorem strength.

Evidence can fail in several non-equivalent ways. A missing artifact is incompleteness. A subject mismatch is invalid composition. A failed mutation shows that the exercised checker does not enforce an intended property. A refinement mismatch shows disagreement between selected executable and reference observations. A verifier disagreement shows that the certificate interpretation is unstable. Fail-closed reporting should preserve these diagnoses, since each has a different repair. Regenerating a filename does not repair a failed fixture, and rerunning a build does not repair a source-identity mismatch.

Archival practice is part of the scholarly claim. An evidence citation should resolve to immutable or content-addressed artifacts containing the subject SHA, run identity, report schema version, fixture identifiers, certificate status, missing-evidence list, and verifier verdict. If artifacts remain only in an expiring CI service, the paper can describe the method but cannot independently support exact historical counts from the checkout alone. A repository release, durable archive, or deposited bundle should therefore accompany any numerical release statement.

The independent verifier should be operationally separate enough to catch generator mistakes, but independence is not absolute. It may share parsers, schemas, libraries, or assumptions with the generator. The assurance case should record that trusted computing base and avoid treating a second JSON verdict as epistemically independent merely because it has a different filename. Mutation of certificate fields, cross-subject fixtures, and intentionally incomplete evidence sets are useful negative cases for the verifier itself.

Finally, release is revocable in the ordinary scientific sense. Discovery of a missing premise, incorrect fixture, compromised toolchain, or misidentified subject does not make the earlier artifact disappear; it changes the current assessment of its force. A superseding record should retain the prior evidence, state the reason for withdrawal, and bind the correction to a new subject. This preserves the audit trail and prevents later readers from confusing historical execution success with a standing warranty of legal or software correctness.

\section{Conclusion and Declarations}
\label{sec:conclusion}

The architecture establishes a conservative composition discipline for legal reasoning. \formalized{} results preserve request identity across normalized observations and graph transitions; retain failures and nonempty partial obligations; compute a finite least Horn closure; constrain structured arguments and attack provenance; implement finite grounded, complete, preferred, and stable semantics; distinguish skeptical from credulous queries; prohibit cross-branch composition; expose solver incompleteness and absent authority; evaluate dimension-indexed quantities exactly; union open obligations; take coordinatewise weaker trust; preserve closure under add-only updates; guard as-of applicability; bind receipts to task, digest, time, and authority; and propagate taint without majority washing.

The overall conservativity principle can be summarized as
\begin{equation}
Assurance(x\otimes y)\preceq Assurance(x)\wedge Assurance(y),
\label{eq:overall-conservativity}
\end{equation}
where $\otimes$ is defined only when request, scope, and branch constraints permit composition. This equation is \derived{} as a cross-layer summary; its component laws are \formalized{}.

Probability, analogy strength, graph similarity, explanation quality, privacy-budget selection, and liability allocation remain \conjectured{} extensions. The formal system can constrain how their evidence is transported and how failures are exposed, but it cannot validate a semantics that has not been supplied.

The result is not a proof of law or a substitute for adjudication. It is a machine-checkable boundary around selected computations. That boundary is useful precisely because it refuses to turn build success into factual truth, formal validity into legal authority, or repeated AI agreement into verification.

\subsection*{Limitations}

ULM16 supplies concrete normal-form composition instances, not closure of every external TheorySpec. The typed graph does not prove a dependent payload refinement for every node kind. Argument coverage is equality to a frozen expected carrier, not universal generator completeness. The exact evaluator proves denotational agreement rather than statutory correctness. Incremental correctness is an extensional contract rather than a proof of a production worklist. Banach convergence is conditional on contraction, completeness, and non-emptiness. Mutation and cross-repository fixtures cover specified cases rather than all behavior.

\subsection*{Funding}

This research received no specific grant from any funding agency in the public, commercial, or not-for-profit sectors.

\subsection*{Declaration of Interests}

The author declares no known competing financial interests or personal relationships that could have appeared to influence the work.

\subsection*{Data and Code Availability}

No new empirical dataset was created. Formal sources, scripts, manifests, and public documentation are available in the Legal Math Modeling repository \citep{LegalMathModeling2026}. Release claims remain subject- and artifact-bound.

\subsection*{Ethics}

The research involved no human participants, personal data, customer matters, or confidential legal records.

\subsection*{CRediT Author Statement}

Laupinco: Conceptualization, Methodology, Formal analysis, Software, Validation, Investigation, Resources, Data curation, Writing---original draft, Writing---review and editing, Project administration.

\subsection*{AI Disclosure}

AI-assisted tools supported literature-search planning, structural drafting, language revision, and formatting. The author directed the work, checked claims against formal artifacts and cited sources, made all legal and scholarly judgments, and accepts full responsibility for the manuscript.

\subsection*{Research Agenda}

The next formal steps are deliberately modular. A probability extension should define measurable carriers, kernels, calibration claims, and executable refinement. An analogy extension should distinguish matching factors, legally decisive differences, authority, and temporal validity. A graph extension should either prove a genuine metric after a justified embedding or retain its score as ranking-only. A privacy extension should bind adjacency, sensitivity, mechanisms, accounting, and authority. An explanation extension should separate structural traceability from human comprehension and causal faithfulness. A liability extension should encode jurisdiction-specific elements without inferring them from technical failure alone.

Each extension should enter through the same sequence: define the semantics, state negative constraints, prove local invariants, build executable fixtures, generate subject-bound evidence, and obtain authorized legal review. This sequence does not guarantee a correct legal system. It does make unsupported transitions visible and gives reviewers a precise point at which to disagree.

The architecture can therefore grow without broadening existing proofs by rhetoric. New theorem families increase the proved surface only when their declarations elaborate and their release evidence is bound to the same source. New runtime components increase executable assurance only when refinement witnesses cover them. New legal conclusions acquire authority only through the applicable human and institutional process.

Evaluation should mirror this layered structure. Formal evaluation asks whether declarations elaborate without prohibited assumptions and whether composition theorems match their stated interfaces. Engineering evaluation asks whether baseline fixtures pass, designed mutations fail, and optimized implementations agree with reference computations. Human-centered evaluation asks whether reports preserve decisive reasons, counterarguments, uncertainty, and actionable recovery information. Legal evaluation asks whether the governing sources, facts, burdens, and institutional authority are correct for the concrete matter. Reporting one layer cannot substitute for reporting the others.

This account favors bounded results over universal slogans. The system does not establish that mathematics settles law. It shows that mathematics can prevent specified transformations from hiding provenance, incompleteness, disagreement, temporal drift, unit errors, weak trust, or missing authority. Those negative guarantees are meaningful infrastructure for computational law because they constrain how uncertainty and institutional judgment move through software without claiming to eliminate either.


\bibliographystyle{IEEEtran}
\bibliography{references}
\end{document}
